% ===================================================================
%  तालिका सं. 9 — पहाड़ और जल-नगरी। — जंगल और शिकार।
%  Traduction 2026 d'apres texto/io, controlee sur texto/fr et
%  texto/en. Consigne : voir texto/hi/10-talika-01.tex.
% ===================================================================

%%K t09-apar-1 apar
\VUsaut{4.0mm}
\VUtitre{15.0pt}{60}{तालिका सं. 9}
\VUsaut{3.0mm}

%%K t09-tit-1 sub
\VUcentre{12.0pt}{0}{\textit{पहला दृश्य।}}
\VUsaut{3.0mm}

%%K t09-tit-2 sub
\VUpk{11.4pt}{40}{पहाड़। --- जल-नगरी।}
\VUsaut{3.0mm}

%%K t09-c1-01-1 p
1. --- मैं एक सप्ताह से प्रात्ज़ में हूँ। पिरेनीज़ की अद्भुत \VUgras{पर्वत-माला}
\textsuperscript{(1)} के सामने खड़े होकर मुझे जो विस्मय हुआ, उसे मैं कह नहीं
सकता; उसकी \VUgras{चोटी-रेखा} \textsuperscript{(2)} नीले आकाश पर किसी विशाल झालर
की तरह उभरती है।

%%K t09-c1-02-1 p
2. --- मैंने कल्पना न की थी कि पिरेनीज़ के \VUgras{शिखर} \textsuperscript{(3)}
इतनी \VUgras{ऊँचाई} तक पहुँच सकते हैं, और मैं उन शानदार \VUgras{हिमनदों}
\textsuperscript{(5)} तथा बर्फ़ से ढकी \VUgras{चोटियों} \textsuperscript{(4)} को
देखकर चकित रह गया, जिनसे इतने भयानक \VUgras{हिमस्खलन} उतरते हैं।

%%K t09-tit-3 sub
\VUsaut{3.0mm}
\VUpk{10.2pt}{40}{पहाड़ी दृश्य।}
\VUsaut{3.0mm}

%%K t09-c1-03-1 p
3. --- आने के अगले ही दिन मैं प्रात्ज़ के पास उस पुराने बुझे \VUgras{ज्वालामुखी}
\textsuperscript{(6)} तक गया। \VUgras{मुख} के नंगे बंजर किनारे पर अब भी वे निशान
साफ़ दिखाई देते हैं जो उस \VUgras{लावा} ने छोड़े थे, जो कई सदियों पहले
\VUgras{पहाड़ की ढाल} \textsuperscript{(7)} से बहा था। एक \VUgras{ढलान} पर मुझे उस
लावे के कुछ टुकड़े मिल ही गए।

%%K t09-c1-04-1 p
4. --- प्रात्ज़ सीमा पर बसा है, और जो \VUgras{किला} \textsuperscript{(8)} फ़्रांस
और स्पेन के बीच के \VUgras{दर्रे} \textsuperscript{(27)} की रखवाली करता है, वह
राइन के किनारे के उन पुराने दुर्गों जैसा ही है जो अपनी चट्टान के ऊपर से किसी
\VUgras{शिकार} की ताक में लगते हैं।

%%K t09-c1-05-1 p
5. --- एक विशाल \VUgras{गरुड़} \textsuperscript{(9)}, जिसका घोंसला
\VUgras{गढ़} \textsuperscript{(8)} से दूर नहीं, प्रायः मेरा ध्यान खींचता है। यह
\VUgras{शिकारी पक्षी}, अपनी बलवान \VUgras{चोंच} के साथ, प्रायः अपने बच्चों के लिए
कोई मेमना या बकरी का बच्चा अपने \VUgras{पंजों} में दबाए लाता है। उसके
\VUgras{पंख} इतने बड़े हैं कि वह अपने भारी बोझ के साथ देर तक तैर सकता है।

%%K t09-tit-4 sub
\VUsaut{3.0mm}
\VUpk{10.2pt}{40}{पैदल चलने वाला।}
\VUsaut{3.0mm}

%%K t09-c1-06-1 p
6. --- यहाँ की सबसे सुखद सैर को के \VUgras{झरने} \textsuperscript{(10)} तक जाना
है। \VUgras{गिरता पानी} बगूला बनकर नीचे आता है और फिर नगर के पश्चिम के
\VUgras{देवदार वन} \textsuperscript{(11)} में एक सुंदर \VUgras{नाले} के रूप में
लुप्त हो जाता है। हम इस वन में से \VUgras{मौसम-केंद्र} \textsuperscript{(12)} तक
चढ़े, और \VUgras{मीनार} के ऊपर से हमने इलाक़े का सारा \VUgras{विहंगम दृश्य}
समेट लिया। \VUgras{प्राकृतिक दृश्य} अत्युत्तम है, और लगता है \VUgras{प्रकृति}
ने इस देश पर अपने सब कोष लुटा दिए हैं।

%%K t09-c1-07-1 p
7. --- नीचे उतरने के लिए हमने \VUgras{रज्जु-रेल} \textsuperscript{(13)} ली और एक
छोटे \VUgras{स्टेशन} पर रुके, एक पुराने \VUgras{सामंती दुर्ग} के
\VUgras{खंडहरों} \textsuperscript{(14)} के पास, जिसमें कभी प्रात्ज़ के सामंत रहते
थे।

%%K t09-c1-08-1 p
8. --- बेचारा दुर्ग! वह क्या सोचता होगा जब अपने नीचे उस सुथरी
\VUgras{जल-नगरी} \textsuperscript{(15)} को देखता है, जो आज एक फैशनेबल
\VUgras{तापीय आरोग्य-स्थल} है?

%%K t09-c1-08-2 p
\VUgras{जलवायु} इतनी सुखद है और \VUgras{खनिज जल} इतना गुणकारी कि
\VUgras{आरोग्यशाला} \textsuperscript{(17)} में लिया गया \VUgras{उपचार-क्रम} सचमुच
सुख की बात है।

%%K t09-tit-5 sub
\VUsaut{3.0mm}
\VUpk{10.2pt}{40}{एक सुहावनी जल-नगरी।}
\VUsaut{3.0mm}

%%K t09-c1-09-1 p
9. --- सचमुच, ठहरने को सुखद बनाने के लिए सब कुछ किया गया है। रेल ले जाने के लिए
एक बड़ा \VUgras{सेतु-पुल} \textsuperscript{(16)} बनाया गया; \VUgras{जल-भवन} का
\VUgras{उपवन} \textsuperscript{(18)}, जहाँ \VUgras{संगीत-सभाएँ} होती हैं, बड़े
सलीके से सजाया गया है। कई सुडौल \VUgras{कोठियाँ} \textsuperscript{(19)}
\VUgras{जुआघर} \textsuperscript{(21)} के चारों ओर बनी हैं, और एक बहुत अच्छी तरह रखी
\VUgras{राजमार्ग} \textsuperscript{(22)} आना-जाना बहुत सहज कर देती है।

%%K t09-c1-10-1 p
10. --- एक सादा \VUgras{बाँध} \textsuperscript{(23)} \VUgras{सड़क}
\textsuperscript{(20)} को असली \VUgras{घाटी} \textsuperscript{(25)} से अलग करता है,
जिसके तल में एक सुडौल छोटी \VUgras{झील} \textsuperscript{(26)} है, जो एक साफ़
\VUgras{नाले} \textsuperscript{(24)} को पानी देती है।

%%K t09-c1-11-1 p
11. --- घाटी के दूसरी ओर लोहे की एक \VUgras{खान} \textsuperscript{(28)} खोदी जा
रही है। कई बार मैं \VUgras{खनिकों} को \VUgras{कुएँ} में उतरते देखने गया, और मैं
उस \VUgras{सुरंग} में भी गया जहाँ वे दिन भर वह \VUgras{कच्ची धातु} निकालते हैं
जिसमें \VUgras{धातु} रहती है।

%%K t09-c1-12-1 p
12. --- खान के पास एक बड़ी \VUgras{ढलाई-शाला} बनाई गई, ताकि निकाली गई संपदा
वहीं काम आ जाए। \VUgras{भट्ठी} \textsuperscript{(29)}, जो यहाँ बहुतायत से मिलने
वाले \VUgras{कोयले} से गरम की जाती है, जल्दी और सस्ते में \VUgras{कच्चा लोहा}
पाने देती है।

%%K t09-c1-13-1 p
13. --- खान से दूर नहीं एक \VUgras{खदान} \textsuperscript{(30)} काटी जा रही है।
बूढ़ा \VUgras{खदान-मज़दूर} अपनी \VUgras{कुदाल} से जो तोड़ता है वह न
\VUgras{ग्रेनाइट} है, न \VUgras{पोर्फ़री}, न \VUgras{बलुआ पत्थर}, बल्कि घर बनाने
का साधारण \VUgras{इमारती पत्थर}।

%%K t09-c1-14-1 p
14. --- इस देश के सबसे सुंदर गहनों में एक है \VUgras{देवदार}
\textsuperscript{(31)}। हर जगह — किसी \VUgras{खड्ड} \textsuperscript{(32)} के तल
में, किसी \VUgras{प्रपात} \textsuperscript{(33)} के किनारे, किसी \VUgras{गह्वर}
\textsuperscript{(34)} या कगार के पास — उसकी पतली सुइयाँ अपना हरा स्वर जोड़ देती
हैं।

%%K t09-tit-6 sub
\VUsaut{3.0mm}
\VUpk{10.2pt}{40}{पैदल चलना।}
\VUsaut{3.0mm}

%%K t09-c1-15-1 p
15. --- मैं अपनी छुट्टियों का लाभ लंबी सैर के लिए उठा रहा हूँ। जो
\VUgras{पगडंडियाँ} \textsuperscript{(35)} पहाड़ पर चक्कर काटती चढ़ती हैं, उन पर
कई \VUgras{किलोमीटर}, और प्रायः कई \VUgras{कोस}, दूरी का पता चले बिना तय हो जाते
हैं।

%%K t09-c1-15-2 p
\VUgras{मील-पत्थर} \textsuperscript{(36)} और \VUgras{दिशा-सूचक}
\textsuperscript{(37)} रास्तों को चिह्नित करते हैं, जिन पर प्रायः कोई जवान
\VUgras{पहाड़ी लड़का} \textsuperscript{(38)} मिलता है, बग़ल में
\VUgras{झोला} \textsuperscript{(40)} लिए और एक \VUgras{मरमोट} \textsuperscript{(39)}
उठाए, या कोई \VUgras{बंजारा} \textsuperscript{(41)}, जिसकी \VUgras{रोएँदार टोपी}
\textsuperscript{(42)} उसके फटे पुराने \VUgras{ओढ़ने} \textsuperscript{(43)} से
अजीब मेल खाती है। वह एक सधे \VUgras{भालू} \textsuperscript{(44)} को
\VUgras{ज़ंजीर} से लिए चलता है।

%%K t09-tit-7 sub
\VUpk{10.2pt}{40}{चढ़ाई।}
\VUsaut{3.0mm}

%%K t09-c1-16-1 p
16. --- लगभग हर दिन मैं एक \VUgras{गाइड} \textsuperscript{(48)} के साथ
\VUgras{चढ़ाई} करने जाता हूँ, और मुझे बड़ा गर्व होता है जब पीठ पर अपना
\VUgras{पिट्ठू} \textsuperscript{(47)} और हाथ में अपनी \VUgras{लाठी} लिए, किसी
सच्चे \VUgras{सैलानी} \textsuperscript{(46)} की तरह, मैं \VUgras{चट्टानों}
\textsuperscript{(45)} पर चढ़ाई करता हूँ।

%%K t09-c1-17-1 p
17. --- पहाड़ की चोटी से मैदान के लोग चींटियों से बड़े नहीं लगते;
\VUgras{भेड़ों का रेवड़} \textsuperscript{(50)}, जो \VUgras{चरागाह}
\textsuperscript{(49)} में चरता है, बच्चों का खिलौना लगता है। एक \VUgras{चीड़}
\textsuperscript{(51)}, या \VUgras{लकड़ी का पहाड़ी घर} \textsuperscript{(52)} तक,
हरे पर एक बिंदु से अधिक नहीं।

%%K t09-c1-18-1 p
18. --- इन सैरों में हमें \VUgras{पगडंडी} \textsuperscript{(53)} पर प्रायः कोई
\VUgras{शामोआ का शिकारी} \textsuperscript{(54)} मिलता है, जो अभी-अभी मारा हुआ
जानवर कंधे पर लिए होता है।

%%K t09-c1-19-1 p
19. --- मैंने अपने पादप-संग्रह में \VUgras{फ़र्न} \textsuperscript{(55)} का एक बहुत
अनोखा पत्ता और \VUgras{हीदर} \textsuperscript{(57)} की एक सुंदर गुलाबी टहनी रखी
है, जो मैंने पिछली सैर में एक छोटी \VUgras{गुफा} \textsuperscript{(56)} के मुँह पर
बटोरी थी।

%%K t09-tit-8 sub
\VUsaut{3.0mm}
\VUcentre{12.0pt}{0}{\textit{दूसरा दृश्य।}}
\VUsaut{3.0mm}

%%K t09-tit-9 sub
\VUpk{11.4pt}{40}{जंगल। --- शिकार।}
\VUsaut{3.0mm}

%%K t09-c2-01-1 p
1. --- कल मैंने जंगल में एक आनंददायक दिन बिताया। वहाँ रहने वाले जानवरों और
\VUgras{कीड़ों} \textsuperscript{(5)} के व्यस्त जीवन ने मुझे बहुत रुचिकर लगा।

%%K t09-c2-01-2 p
\VUgras{मकड़ी} \textsuperscript{(1)}, जो दो डालों के बीच अपना \VUgras{जाला}
\textsuperscript{(2)} बुनकर गुज़रती \VUgras{मक्खी} \textsuperscript{(3)} पकड़ती है,
\VUgras{घोंघा} \textsuperscript{(4)}, जो अपना ख़ोल लिए मुश्किल से आगे बढ़ता है,
वह भृंग जो अपने शिकार की टोह में है, वह मेहनती चींटी जो \VUgras{बाँबी}
\textsuperscript{(6)} के पास काम में जुटी है — ये सब नन्हे जीव असाधारण उत्साह से
आते-जाते और काम करते रहे।

%%K t09-c2-02-1 p
2. --- एक \VUgras{साँप} \textsuperscript{(7)}, जिसे पहले मैंने \VUgras{नाग} समझा
पर जो केवल एक निर्दोष \VUgras{पनिहा} था, एक तने के नीचे कुंडली मारे था, और
बीच-बीच में अपना पतला सिर निकालता था, जो सदा डँसने को तैयार लगता था। मेरे सामने
एक बड़ा \VUgras{स्लग} \textsuperscript{(8)} भारी चाल से रेंग रहा था, उन स्वादिष्ट
नन्ही \VUgras{जंगली स्ट्रॉबेरियों} \textsuperscript{(11)} में से एक चट कर लेने के
बाद, जो यहाँ बहुतायत से उगती हैं।

%%K t09-c2-03-1 p
3. --- ज़मीन \VUgras{जंगली फूलों} से बिछी थी: \VUgras{प्रिमरोज़}
\textsuperscript{(9)}, \VUgras{पेरिविंकल} \textsuperscript{(10)} और और भी बहुत से
पौधे, जिन्हें मैं जानता न था, \VUgras{घास के गुच्छों} \textsuperscript{(12)} के
बीच खिले थे।

%%K t09-c2-04-1 p
4. --- इस इलाक़े में \VUgras{ट्रफ़ल} लगभग उतना ही आम है जितना \VUgras{कुकुरमुत्ता}
\textsuperscript{(14)}, और एक वनरक्षक का \VUgras{सूअर का बच्चा} \textsuperscript{(13)}
लालच से खोज में लगा था कि कुछ मिल जाए।

%%K t09-tit-10 sub
\VUsaut{3.0mm}
\VUpk{10.2pt}{40}{चोरी का शिकार।}
\VUsaut{3.0mm}

%%K t09-c2-05-1 p
5. --- मैंने एक दृश्य का \VUgras{अंत} देखा जिसने मुझे बहुत हँसाया। अपने
\VUgras{बासे कुत्ते} \textsuperscript{(15)} की नाक की बदौलत \VUgras{वनरक्षक}
\textsuperscript{(16)} ने अभी-अभी एक \VUgras{चोर-शिकारी} \textsuperscript{(22)} को
उस समय रँगे हाथ पकड़ा था जब वह अपना मारा \VUgras{शिकार} \textsuperscript{(17)}
उठा ले जाने की तैयारी में था। पकड़े जाने से अच्छा उसे अपना शिकार छोड़ देना लगा,
सो चोर-शिकारी भाग गया, और \VUgras{ख़रगोश} \textsuperscript{(18)},
\VUgras{हिरनी} \textsuperscript{(19)}, \VUgras{बारहसिंगा} \textsuperscript{(20)} और
\VUgras{जंगली सूअर} \textsuperscript{(21)} घास पर पड़े रह गए, वनरक्षक की ख़ुशी के
लिए।

%%K t09-c2-06-1 p
6. --- जंगल में पेड़ों की जितनी क़िस्में हैं, वह चकित कर देती है।
\VUgras{बीच} \textsuperscript{(23)} और \VUgras{भोजपत्र} \textsuperscript{(26)},
\VUgras{चीड़}, जिनसे \VUgras{राल} \textsuperscript{(27)} मिलती है,
\VUgras{बलूत} \textsuperscript{(29)} और और भी बहुत से अपनी डालें मिलाते हैं।

%%K t09-c2-06-2 p
ऊँचे पेड़ों के तनों से लिपटी \VUgras{आइवी} \textsuperscript{(24)}
\VUgras{ऊँचे वन} \textsuperscript{(25)} को एक चमकीला हरा रंग देती है, जो उस
\VUgras{काई} \textsuperscript{(31)} की फीकी कोमल छाया से सुखद ढंग से मिल जाता है
जिसमें \VUgras{हरा कठफोड़वा} \textsuperscript{(30)} कीड़े खोजता है।

%%K t09-tit-11 sub
\VUsaut{3.0mm}
\VUpk{10.2pt}{40}{जंगल का दृश्य।}
\VUsaut{3.0mm}

%%K t09-c2-07-1 p
7. --- पुराने बलूत मुझे मोह लेते हैं। उनकी बड़ी बल खाती \VUgras{जड़ें}
\textsuperscript{(32)}, आधी ज़मीन से बाहर, उन्हें बल और ओज का शानदार रूप देती
हैं, और मुझे अचरज होता है कि \VUgras{मालिक} \textsuperscript{(35)} उन्हें कटवाकर उस
\VUgras{आरी}
\textsuperscript{(34)} को सौंपने को कैसे राज़ी हो जाता है जो
\VUgras{लकड़हारे} \textsuperscript{(33)} की है। बेचारे \VUgras{तने} \textsuperscript{(36)},
अपनी \VUgras{छाल} \textsuperscript{(37)} से छीले या \VUgras{कुल्हाड़े}
\textsuperscript{(38)} से चीरे, जो \VUgras{दिहाड़ी मज़दूर} \textsuperscript{(39)}
का है, मुझे दुख देते हैं।

%%K t09-c2-08-1 p
8. --- पास ही एक बूढ़े \VUgras{कोयला बनाने वाले} \textsuperscript{(41)} ने अपने
\VUgras{बेलचे} से कोयले का एक \VUgras{ढेर} \textsuperscript{(42)} लगा रखा था, और
एक बुढ़िया कटे पेड़ों की टहनियों से बना \VUgras{सूखी लकड़ी} का एक \VUgras{गट्ठर}
\textsuperscript{(40)} ले जा रही थी।

%%K t09-tit-12 sub
\VUsaut{3.0mm}
\VUpk{10.2pt}{40}{शिकार।}
\VUsaut{3.0mm}

%%K t09-c2-09-1 p
9. --- मैं \VUgras{वनरक्षक के घर} \textsuperscript{(46)} में थोड़ा सुस्ताने जाने
वाला था, पर उस छोटे \VUgras{ताल} \textsuperscript{(43)} तक पहुँचकर, जिस पर एक
सुंदर \VUgras{हंस} \textsuperscript{(45)} तैर रहा था, मैंने देखा कि हाल की एक
\VUgras{बाढ़} \textsuperscript{(44)} ने रास्ते को सचमुच का \VUgras{दलदल} बना दिया
है। मुझे चक्कर काटना पड़ा, और \VUgras{देवदार} \textsuperscript{(49)},
\VUgras{चिनारों} \textsuperscript{(48)} तथा \VUgras{एल्म} \textsuperscript{(47)} के
पास से गुज़रते हुए, जो जंगल के किनारे खड़े हैं, मैंने देखा कि एक
\VUgras{झाड़ी} \textsuperscript{(54)} से निकलकर एक शानदार \VUgras{बारहसिंगा}
\textsuperscript{(50)} भागा, जिसके पीछे एक न थकने वाला \VUgras{शिकारी कुत्तों का
झुंड} \textsuperscript{(51)} लगा था, जिसे \VUgras{नरसिंगा} \textsuperscript{(53)} उकसा रहा था,
\VUgras{घोड़े पर सवार शिकारियों} \textsuperscript{(52)} का।
बेचारा जानवर बिलकुल थक चुका था। वह मुझसे कुछ ही क़दम पर मरता हुआ गिर पड़ा।

%%K t09-c2-10-1 p
10. --- वनरक्षक का बेटा, जिसे \VUgras{शिकार} के शोर ने बाहर खींच लिया, घर से
निकला, और हम दोनों फिर जंगल में लौटे। उसने एक पुराने बलूत की डाल पर एक
\VUgras{नीलकंठ} \textsuperscript{(56)} देखा था। उसे लगा कि उसका घोंसला
\VUgras{पत्तों} \textsuperscript{(58)} में छिपा होगा, और वह उसे लेना चाहता था।

%%K t09-c2-10-2 p
\VUgras{गिलहरी} \textsuperscript{(61)} की तरह फुरतीला, वह एक \VUgras{डाल}
\textsuperscript{(55)} से दूसरी डाल चढ़ता गया, पर उसे कुछ न मिला और वह केवल एक बूढ़े
\VUgras{कौवे} \textsuperscript{(60)} को चौंका सका, जो एक \VUgras{शाखा}
\textsuperscript{(57)} पर बैठा था।
\VUsaut{4.0mm}
\VUfilet{22mm}
